% !TEX program = lualatex
\DocumentMetadata{pdfstandard=UA-2,pdfversion=2.0,lang=en-US,tagging=on}
\documentclass{./../../latex/tests}
\begin{document}
\thispagestyle{plain}
\myheader{Midterm Sample 2}
\rhead{Midterm Sample 2}

\vspace{0.5em}
\testHeader{65}{18}

\begin{enumerate}
\item (5 pts) Answer the following questions. 
\begin{enumerate}
\item (1 pt) Consider a mapping $f(x)$. For two distinct values of $x$, $x_0$ and $x_1$, $f(x_0) = f(x_1)$. Is $f$ a valid function? Answer yes or no. \\~\\
\item (2 pts) Find the union and intersection for the following sets:
\[ A = \{x: x \text{ is an even number} \} \quad \quad B = \{2, 4, 8 \} \] \\~\\
\item (1 pt) Consider the following two-variable function: 
  \[ f(x,y) = x+y \]
  where $x \in (0,1)$ and $y \in (0,1)$. What is the range of $f$? \\~\\
 
\item (1 pt) Given a system of linear equations $A x = b$, if $|A|=5$, what can we say about the solution for this system of equations?
\begin{itemize}
	\item[$\square$] Has no solution. 
	\item[$\square$] Has a unique solution.
	\item[$\square$] Has infinitely many solutions.
	\item[$\square$] None of the above \\~\\
\end{itemize}  
\end{enumerate} 

\newpage
\item (5 pts) Consider the following matrix \[A = I - X(X^T X)^{-1}X^T \]
\begin{enumerate}
\item (3 pts)  Is $A$ a square matrix? Show your work or reasoning that led you to this conclusion. \vspace{2.5cm}
\item (2 pts) Prove that $A$ is idempotent i.e. $A A = A$. \vspace{5cm}
\end{enumerate} 

\item (8 pts) Consider the following system of equations:
\begin{align*}
x-2z & = 2 \\
y+z & = 12 \\
x+y+z &= 24	
\end{align*}
\begin{enumerate}
  \item (1 pt) Write this system of equations in matrix format i.e., \[ Av=b \]
  What is $A$, $v$, and $b$ equal to?
  \newpage
   \item (2 pts) Calculate the adjoint of $A$. \vspace{8.5cm}
  \item (2 pts) Calculate the determinant of $A$. Is $A$ nonsingular? \vspace{4.5cm}
  \item (1 pt) If you premultiply $A^{-1}$ on both sides of the equation $ Av=b $, you should be able to derive an expression to solve for $v$. Write down this expression. \vspace{2cm}
  \item (2 pts) Using the expression in $(d)$ solve for $v^*$. 
\end{enumerate}

\newpage
\end{enumerate}

\end{document}